\documentclass[12pt,a4paper,oneside,openany]{report}
\usepackage[utf8]{inputenc}
\usepackage{times}
\usepackage{geometry}
\geometry{a4paper,top=25mm,bottom=20mm,left=30mm,right=20mm}
\usepackage{setspace}
\onehalfspacing
\usepackage{parskip}
\usepackage{amsmath,amssymb,amsthm}
\usepackage{graphicx}
\graphicspath{{figures/}}
\usepackage{float}
\usepackage{booktabs,longtable,array,multicol,caption}
\usepackage{xcolor}
\usepackage{listings}
\lstset{language=Python,basicstyle=\ttfamily\small,
  keywordstyle=\color{blue},commentstyle=\color{gray},
  stringstyle=\color{orange},breaklines=true,frame=single,
  numbers=left,numberstyle=\tiny\color{gray},xleftmargin=15pt,
  extendedchars=true,inputencoding=utf8}
\usepackage{fancyhdr}
\pagestyle{fancy}
\fancyhf{}
\fancyfoot[C]{\thepage}
\def\headrulewidth{0pt}
\def\footrulewidth{0pt}
\fancypagestyle{plain}{\fancyhf{}\fancyfoot[C]{\thepage}\def\headrulewidth{0pt}}
\usepackage{hyperref}
\hypersetup{colorlinks=true,linkcolor=black}

\begin{document}

% ================================================================
% TITLE PAGE
% ================================================================
\begin{titlepage}
\centering
\vspace*{3cm}

{\Large\textbf{UNDERSTANDING FIGURES}}

\vspace{2cm}

{\normalsize A comprehensive guide to the figures in}
\\
{\normalsize ``Quantifying Privacy Risk in Token-Based CBDC: A}
\\
{\normalsize Normalised Metadata De-Anonymization Framework for Nepal''}

\vspace{3cm}

\end{titlepage}

% ================================================================
% ANALYSIS 3 - SIX-SCHEME COMPARISON
% ================================================================
\chapter{Analysis 3: Six-Scheme Comparison}
\label{chap:a3}

This chapter explains Figure 4.1 (a3\_comparison.png) and Figure 4.2 (a3\_lambda.png).

% ----------------------------------------------------------------
\section{Figure 4.1a: R\_norm Bar Chart}
\label{sec:f4.1a}

\subsection{What Type of Figure This Is}
This is a **vertical bar chart** (also known as a column chart). It displays the normalised privacy risk score ($R_{\mathrm{norm}}$) for six different cryptographic schemes used in the CBDC system.

\subsection{What Is Being Shown}
The x-axis shows the six cryptographic schemes:
\begin{itemize}
    \item Plain Hash Token
    \item Blind Signature
    \item Hybrid Scheme
    \item Ring Signature
    \item Stealth Address
    \item ZKP Token
\end{itemize}

The y-axis shows the $R_{\mathrm{norm}}$ value ranging from 0 to 0.4 (representing the normalized risk score between 0 and 1).

Each bar is colored differently to distinguish the schemes:
\begin{itemize}
    \item Red: Plain Hash Token (highest risk)
    \item Orange: Blind Signature
    \item Yellow: Hybrid Scheme
    \item Green: Ring Signature
    \item Cyan: Stealth Address
    \item Blue: ZKP Token (lowest risk)
\end{itemize}

\subsection{Why This Graph Type Was Chosen}
A bar chart is the ideal choice here because:
\begin{enumerate}
    \item It allows easy comparison between categorical variables (the six schemes)
    \item The length of each bar directly corresponds to the risk value, making visual comparison intuitive
    \item It clearly shows the ranking from highest to lowest risk
    \item The color coding adds an additional visual cue for quick identification
\end{enumerate}

\subsection{Code Explanation}
The following Python code generates this figure:

\begin{lstlisting}[language=Python]
# Create figure with 3 subplots
fig, axes = plt.subplots(1, 3, figsize=(18, 7))

# Subplot 1: R_norm bar chart
bars = axes[0].bar(
    range(len(SCHEMES)), R_norm_vals,
    color=[COLORS[s] for s in SCHEMES],
    edgecolor=EDGE, linewidth=1.2, width=0.6)
axes[0].set_xticks(range(len(SCHEMES)))
axes[0].set_xticklabels(LABELS, color="black", fontsize=10)
axes[0].set_ylabel("R_norm  [0, 1]", fontsize=12)
axes[0].set_title("R_norm Comparison", fontsize=12, fontweight="bold")
axes[0].set_ylim(0, 0.4)
\end{lstlisting}

\textbf{Line-by-line explanation:}
\begin{itemize}
    \item \texttt{fig, axes = plt.subplots(1, 3, figsize=(18, 7))} - Creates a figure with 1 row and 3 columns of subplots, with size 18x7 inches
    \item \texttt{bars = axes[0].bar(...)} - Creates a bar chart in the first subplot (axes[0])
    \item \texttt{range(len(SCHEMES))} - Creates x-positions for each scheme (0, 1, 2, 3, 4, 5)
    \item \texttt{R\_norm\_vals} - The actual risk values to plot
    \item \texttt{color=[COLORS[s] for s in SCHEMES]} - Assigns colors to each bar based on the COLORS dictionary
    \item \texttt{edgecolor=EDGE} - Sets the border color of each bar
    \item \texttt{axes[0].set\_xticks(...)} - Sets the x-axis tick positions
    \item \texttt{axes[0].set\_xticklabels(...)} - Sets the labels for each tick (scheme names)
    \item \texttt{axes[0].set\_ylabel(...)} - Sets the y-axis label
    \item \texttt{axes[0].set\_ylim(0, 0.4)} - Sets the y-axis limits from 0 to 0.4
\end{itemize}

\subsection{What This Graph Signifies}
This figure demonstrates that:
\begin{enumerate}
    \item \textbf{Plain Hash Token has the highest risk} ($R_{\mathrm{norm}} = 0.317$) - This is expected because it provides no metadata privacy; all attributes are fully exposed.
    \item \textbf{ZKP Token has the lowest risk} ($R_{\mathrm{norm}} = 0.035$) - Zero-knowledge proofs provide the strongest privacy guarantees.
    \item \textbf{There is a clear ranking} - The schemes form a non-overlapping risk ladder from highest to lowest.
    \item \textbf{The difference between best and worst is ~9x} - ZKP Token is approximately 9 times safer than Plain Hash Token.
\end{enumerate}

\subsection{Possible Exam Questions}
\begin{enumerate}
    \item \textbf{Question:} Why does Plain Hash Token have the highest $R_{\mathrm{norm}}$ value?
    \item \textbf{Answer:} Because it provides no cryptographic privacy protection. All metadata attributes (amount, timestamp, merchant, wallet, etc.) are recorded exactly as they are, making users easily re-identifiable.

    \item \textbf{Question:} What does the bar chart tell us about the effectiveness of different cryptographic schemes?
    \item \textbf{Answer:} It shows that scheme choice has a dramatic impact on privacy risk. The 9x difference between Plain Hash Token and ZKP Token demonstrates that cryptographic design is the primary determinant of metadata privacy.

    \item \textbf{Question:} If you were advising NRB on CBDC deployment, which scheme would you recommend based on this figure?
    \item \textbf{Answer:} ZKP Token provides the best privacy protection. However, if computational resources are limited, Ring Signature or Hybrid Scheme (both around 0.06-0.07) are reasonable alternatives.
\end{enumerate}

% ----------------------------------------------------------------
\section{Figure 4.1b: Attribute Contribution Stacked Bar Chart}
\label{sec:f4.1b}

\subsection{What Type of Figure This Is}
This is a **stacked bar chart** showing how each metadata attribute contributes to the overall $R_{\mathrm{norm}}$ score for each cryptographic scheme.

\subsection{What Is Being Shown}
The x-axis shows the six cryptographic schemes (same as Figure 4.1a).

The y-axis shows the cumulative $s(a_i)$ values - the single-attribute uniqueness scores added together.

Each bar is divided into segments, with each segment representing a different attribute:
\begin{itemize}
    \item amount (blue)
    \item timestamp (orange)
    \item merchant (green)
    \item wallet (red)
    \item frequency (purple)
    \item payment\_channel (cyan - Nepal-specific)
    \item transaction\_zone (yellow - Nepal-specific)
\end{itemize}

The height of each bar shows the total contribution from all attributes.

\subsection{Why This Graph Type Was Chosen}
A stacked bar chart was chosen because:
\begin{enumerate}
    \item It shows both the total risk \textbf{and} the composition of that risk
    \item It allows comparison of which attributes contribute most to risk for each scheme
    \item It clearly shows whether risk is concentrated in one attribute or distributed across multiple
\end{enumerate}

\subsection{What This Graph Signifies}
Key observations:
\begin{enumerate}
    \item \textbf{Wallet dominates} - For most schemes, the wallet attribute (red segment) contributes the most to total risk, often being the largest segment in the stack.
    \item \textbf{Scheme affects attribute contribution} - Different schemes "hide" different attributes. For example, Blind Signature reduces the timestamp segment to near-zero.
    \item \textbf{Nepal-specific attributes show zero contribution} - The payment\_channel and transaction\_zone segments are zero for all schemes because these attributes have $s(a_i) = 0$ due to Nepal's demographic clustering.
\end{enumerate}

% ----------------------------------------------------------------
\section{Figure 4.1c: Nepal-Specific Attributes}
\label{sec:f4.1c}

\subsection{What Type of Figure This Is}
This is a simple bar chart showing the $s(a_i)$ values for Nepal-specific attributes only: payment\_channel and transaction\_zone.

\subsection{What It Shows}
The figure displays a single dotted horizontal line at $y = 0$ for all six schemes. This represents that both payment\_channel and transaction\_zone have $s(a_i) = 0$ (zero uniqueness) for all cryptographic schemes.

\subsection{Why It Looks Like This}
This figure appears as a flat line at zero because:

\textbf{The underlying data:}
\begin{itemize}
    \item $s(\text{payment\_channel}) = 0.000$ for all schemes
    \item $s(\text{transaction\_zone}) = 0.000$ for all schemes
\end{itemize}

\textbf{Why these values are zero:}

In Nepal's specific context:
\begin{enumerate}
    \item \textbf{Payment Channel Distribution:} 60\% mobile wallet, 35\% QR code, 5\% other - most users share the same payment channel values, so no one is unique on this attribute.
    \item \textbf{Transaction Zone Distribution:} 35\% Kathmandu Valley, 25\% Terai Urban, 20\% Hill Urban, 12\% Terai Rural, 8\% Hill Rural - the population is heavily concentrated in certain zones, making uniqueness on this attribute essentially zero.
\end{enumerate}

The combination of these skewed distributions means that when calculating $s(a_i)$ (the fraction of users who are unique on attribute $a_i$), the result is zero because most users share the same values.

\textbf{Code Logic:}
\begin{lstlisting}[language=Python]
# In the generate_attributes function:
if attribute == "payment_channel":
    # 60% mobile wallet, 35% QR code, 5% other
    return np.random.choice([0, 1, 2], N, p=[0.60, 0.35, 0.05])

if attribute == "transaction_zone":
    # 35% Kathmandu, 25% Terai Urban, 20% Hill Urban, etc.
    return np.random.choice([0, 1, 2, 3, 4], N, p=[0.35, 0.25, 0.20, 0.12, 0.08])
\end{lstlisting}

Because these probability distributions are heavily skewed, the probability of any single user being unique (having a value no one else has) is effectively zero.

\textbf{Is it correctly formed?} Yes. The dotted horizontal line at $y=0$ correctly represents that there is zero contribution from Nepal-specific attributes to the overall risk score. This is actually a \textit{good thing} - it means Nepal's demographic structure provides passive privacy protection on these attributes.

\textbf{The dotted line style} is used to indicate that there is no meaningful data to display (the bars would all be at height 0), while still making it clear that this was intentionally calculated.

% ----------------------------------------------------------------
\section{Figure 4.2: Lambda Heatmaps}
\label{sec:f4.2}

\subsection{What Type of Figure This Is}
This is a **heatmap matrix** showing the pairwise interaction coefficient $\lambda_{ij}$ for each pair of attributes, for each of the six cryptographic schemes.

A heatmap uses color intensity to represent numerical values - darker colors indicate higher values.

\subsection{What Is Being Shown}
The figure contains 6 subplots (2 rows × 3 columns), one for each cryptographic scheme.

Each subplot is a 7×7 matrix where:
\begin{itemize}
    \item Rows and columns represent the 7 metadata attributes
    \item Each cell shows the $\lambda_{ij}$ value (interaction coefficient) between that pair of attributes
    \item Color intensity represents the value: darker red/orange = higher interaction, lighter = lower
\end{itemize}

The diagonal cells (where row = column) are empty because $\lambda_{ij}$ is not defined for the same attribute pair.

The color scale (legend) on the right shows the values, ranging from 0 to 100 (capped at 100 for display).

\subsection{Why This Graph Type Was Chosen}
A heatmap is ideal here because:
\begin{enumerate}
    \item It can display a large matrix of values (7×7 = 49 values) in a compact, readable format
    \item Color allows quick visual identification of which attribute pairs have the highest interaction
    \item It enables comparison across all 6 schemes in a single figure
    \item The matrix format directly maps to the mathematical $\lambda_{ij}$ concept
\end{enumerate}

\subsection{Code Explanation}
\begin{lstlisting}[language=Python]
# Create 2x3 subplot grid
fig, axes = plt.subplots(2, 3, figsize=(20, 14))

for ax, scheme in zip(axes.flatten(), SCHEMES):
    # Create empty matrix for this scheme
    mat = np.zeros((len(ATTRS_7), len(ATTRS_7)))
    
    # Fill matrix with lambda values
    for i, ai in enumerate(ATTRS_7):
        for j, aj in enumerate(ATTRS_7):
            if i != j:
                mat[i][j] = min(res3[scheme]["lmat"].get(ai, {}).get(aj, 0.0), 100)
    
    # Create heatmap
    sns.heatmap(mat, ax=ax,
        xticklabels=ATTRS_7, yticklabels=ATTRS_7,
        annot=True, fmt=".1f", cmap="YlOrRd",
        annot_kws={"size": 9})
\end{lstlisting}

\textbf{Key lines explained:}
\begin{itemize}
    \item \texttt{mat = np.zeros((7, 7))} - Creates a 7×7 matrix initialized to zeros
    \item \texttt{for i, ai in enumerate(ATTRS\_7)} - Iterates through all attribute pairs
    \item \texttt{mat[i][j] = min(..., 100)} - Stores the lambda value, capped at 100 for display
    \item \texttt{sns.heatmap(...)} - Creates the heatmap with color scale "YlOrRd" (Yellow-Orange-Red)
    \item \texttt{annot=True, fmt=".1f"} - Shows numerical values in each cell with 1 decimal place
\end{itemize}

\subsection{What This Graph Signifies}
\begin{enumerate}
    \item \textbf{Wallet-timestamp and wallet-amount interactions are often high} - These pairs typically show the darkest colors, indicating strong interaction effects
    \item \textbf{ZKP Token shows minimal interaction} - The heatmap is much lighter because ZKP reduces all attribute uniqueness to near-zero
    \item \textbf{Plain Hash Token shows diverse interactions} - Multiple attribute pairs show high lambda values
    \item \textbf{Diagonal is empty} - Lambda is only defined for pairs (i ≠ j), so diagonal cells are left blank
\end{enumerate}

\chapter{Analysis 5: Population Scalability}
\label{chap:a5}

This chapter explains Figure 4.3 (a5\_scalability.png).

% ================================================================
\section{Figure 4.3: Population Scalability}
\label{sec:f4.3}

\subsection{What Type of Figure This Is}
This is a **dual-panel figure**:
\begin{itemize}
    \item Left panel: Line graph showing how $R_{\mathrm{norm}}$ changes as population size $N$ increases
    \item Right panel: Bar chart comparing risk ratios relative to ZKP Token at $N=100,000$
\end{itemize}

\subsection{What Is Being Shown}

\textbf{Left Panel (Line Graph):}
\begin{itemize}
    \item X-axis: Population size $N$ (log scale: 100, 1,000, 10,000, 100,000)
    \item Y-axis: $R_{\mathrm{norm}}$ (0 to 0.6)
    \item 6 lines, one for each cryptographic scheme
    \item X-axis is logarithmic (each tick represents a 10× increase in population)
\end{itemize}

\textbf{Right Panel (Bar Chart):}
\begin{itemize}
    \item X-axis: The six cryptographic schemes
    \item Y-axis: Risk relative to ZKP Token (how many times riskier than ZKP)
    \item At $N=100,000$ (national scale)
\end{itemize}

A vertical dashed line at $N=50$ indicates the VDC (village) pilot scale.

\subsection{Why This Graph Type Was Chosen}
\begin{enumerate}
    \item \textbf{Log scale on x-axis} is necessary because population ranges from 100 to 100,000 (3 orders of magnitude)
    \item \textbf{Dual panel} shows two different but related insights: how risk changes with scale (left), and relative comparison at national scale (right)
    \item \textbf{Line graph} best shows trends over a continuous variable (population size)
    \item \textbf{Bar chart} enables direct comparison of the final ratios
\end{enumerate}

\subsection{What This Graph Signifies}

Key findings from this figure:

\begin{enumerate}
    \item \textbf{Plain Hash Token does not improve with scale} - The red line remains high even at $N=100,000$, showing that population size cannot substitute for cryptographic design
    \item \textbf{ZKP Token is scale-independent} - The blue line is flat at 0.035 across all population sizes
    \item \textbf{Ring Signature and Hybrid Scheme plateau early} - Both reach their minimum around $N=1,000$ and stay flat
    \item \textbf{Plain Hash at N=100,000 is 5.5× riskier than ZKP} - The right panel shows this ratio explicitly
    \item \textbf{Micro-scale (N=50) is shown} - The vertical line indicates village pilot scale, where all schemes except ZKP have very high risk
\end{enumerate}

\chapter{Analysis 7: Mitigation Effectiveness}
\label{chap:a7}

This chapter explains Figure 4.4 (a7\_mitigation.png).

% ================================================================
\section{Figure 4.4: Mitigation Effectiveness}
\label{sec:f4.4}

\subsection{What Type of Figure This Is}
This is a **dual-panel grouped bar chart**:
\begin{itemize}
    \item Left: Before and after mitigation comparison
    \item Right: Percentage reduction achieved by each mitigation strategy
\end{itemize}

\subsection{What Is Being Shown}

\textbf{Left Panel (Before/After):}
\begin{itemize}
    \item X-axis: The six schemes
    \item Y-axis: $R_{\mathrm{norm}}$ (0 to 0.35)
    \item Two bars per scheme: dark (before mitigation), light (after full mitigation)
    \item Shows the dramatic reduction in risk when all mitigations are applied
\end{itemize}

\textbf{Right Panel (Reduction by Strategy):}
\begin{itemize}
    \item X-axis: Five mitigation strategies
    \item Y-axis: Percentage reduction in $R_{\mathrm{norm}}$
    \item Shows how much each individual strategy contributes to the total reduction
\end{itemize}

The five mitigation strategies are:
\begin{enumerate}
    \item Amount bucketing (grouping transaction amounts into ranges)
    \item Timestamp coarsening (reducing time precision to day-of-week only)
    \item Merchant categorisation (grouping merchants into sectors)
    \item Wallet unlinkability (breaking the link between consecutive transactions)
    \item Frequency smoothing (adding random noise to transaction frequency)
\end{enumerate}

\subsection{What This Graph Signifies}
\begin{enumerate}
    \item \textbf{Full mitigation reduces risk by 91-97\%} for all schemes
    \item \textbf{Wallet unlinkability is the most effective single strategy} - shown in the right panel as the highest bar
    \item \textbf{Ring Signature with full mitigation achieves near-ZKP levels} - both reach approximately 0.002
    \item \textbf{No scheme is immune to mitigation} - all benefit significantly
    \item \textbf{Plain Hash Token still has residual risk} - even with mitigation, it has $R_{\mathrm{norm}} = 0.017$
\end{enumerate}

\chapter{Analysis 8: m-Attribute Sensitivity}
\label{chap:a8}

This chapter explains Figure 4.5 (a8\_m\_sensitivity.png).

% ================================================================
\section{Figure 4.5: m-Attribute Sensitivity}
\label{sec:f4.5}

\subsection{What Type of Figure This Is}
This is a **triple-panel figure**:
\begin{enumerate}
    \item Left: Line graph showing $R_{\mathrm{norm}}$ as attributes are added one by one
    \item Center: Line graph showing $R_{\mathrm{max}}$ as attributes are added
    \item Right: Bar chart showing the marginal contribution of each new attribute
\end{enumerate}

\subsection{What Is Being Shown}

\textbf{Left Panel ($R_{\mathrm{norm}}$ vs m):}
\begin{itemize}
    \item X-axis: Number of attributes (m = 1 to 7)
    \item Y-axis: $R_{\mathrm{norm}}$
    \item Shows how risk changes as more attributes are modelled
    \item The vertical dashed line at m=5 separates standard attributes (1-5) from Nepal-specific attributes (6-7)
\end{itemize}

\textbf{Center Panel ($R_{\mathrm{max}}$ vs m):}
\begin{itemize}
    \item X-axis: Number of attributes (m = 1 to 7)
    \item Y-axis: $R_{\mathrm{max}}$ (theoretical maximum)
    \item Shows formula: $R_{\mathrm{max}}(m) = m + \binom{m}{2}$
    \item Values shown: m=1→1, m=2→3, m=3→6, m=4→10, m=5→15, m=6→21, m=7→28
\end{itemize}

\textbf{Right Panel (Delta $R_{\mathrm{norm}}$):}
\begin{itemize}
    \item X-axis: Attributes (added in order)
    \item Y-axis: Change in $R_{\mathrm{norm}}$ when that attribute is added
    \item Shows the "cost" of each additional attribute in terms of privacy risk
\end{itemize}

\subsection{What This Graph Signifies}
\begin{enumerate}
    \item \textbf{Adding Nepal-specific attributes actually DECREASES $R_{\mathrm{norm}}$} - When payment\_channel (m=6) and transaction\_zone (m=7) are added, $R_{\mathrm{norm}}$ goes down because $R_{\mathrm{max}}$ grows faster than $R_{\mathrm{raw}}$
    \item \textbf{Wallet is the highest-risk attribute} - Adding it first causes the largest jump in $R_{\mathrm{norm}}$
    \item \textbf{$R_{\mathrm{max}}$ grows quadratically} - The center panel shows the formula $m + \binom{m}{2}$ in action
    \item \textbf{Nepal's demographic structure provides passive protection} - The counterintuitive decrease proves that Nepal's skewed distributions act as natural privacy
\end{enumerate}

\chapter{Analysis 9: Seed Sensitivity}
\label{chap:a9}

This chapter explains Figure 4.6 (a9\_seed\_sensitivity.png).

% ================================================================
\section{Figure 4.6: Seed Sensitivity}
\label{sec:f4.6}

\subsection{What Type of Figure This Is}
This is a **dual-panel figure**:
\begin{enumerate}
    \item Left: Line graph showing mean $R_{\mathrm{norm}}$ with error bars (2σ) across 30 random seeds
    \item Right: Violin plot showing the distribution of $R_{\mathrm{norm}}$ values for each scheme
\end{itemize}

\subsection{What Is Being Shown}

\textbf{Left Panel (Mean ± 2σ):}
\begin{itemize}
    \item X-axis: The six schemes
    \item Y-axis: $R_{\mathrm{norm}}$ (0 to 0.4)
    \item Points show mean value across 30 runs
    \item Error bars show ±2 standard deviations (covering ~95\% of runs)
    \item Error bars on some schemes are barely visible (very low variance)
\end{itemize}

\textbf{Right Panel (Violin Plot):}
\begin{itemize}
    \item X-axis: The six schemes  
    \item Y-axis: Distribution of $R_{\mathrm{norm}}$ values
    \item The "violin" shape shows the density distribution
    \item White dot shows median, yellow dot shows mean
\end{itemize}

Key statistics shown:
\begin{itemize}
    \item Plain Hash Token: mean=0.323, CV=7.8\%
    \item Blind Signature: mean=0.175, CV=0.3\%
    \item Hybrid Scheme: mean=0.070, CV=21.0\%
    \item Ring Signature: mean=0.069, CV=9.4\%
    \item Stealth Address: mean=0.107, CV=0.1\%
    \item ZKP Token: mean=0.035, CV=0.3\%
\end{itemize}

\subsection{What This Graph Signifies}
\begin{enumerate}
    \item \textbf{The scheme ranking is robust} - No seed produces a ranking reversal; ZKP is always lowest, Plain Hash always highest
    \item \textbf{Hybrid Scheme has highest variance} (CV=21\%) - The wide violin shape shows it is sensitive to random seed
    \item \textbf{ZKP, Blind Signature, and Stealth Address are deterministic} - Their tight error bars and narrow violins show almost no variation
    \item \textbf{Conclusions are not sensitive to random seed choice} - This validates the reproducibility of the analysis
\end{enumerate}

\chapter{Analysis 10: Micro-Scale Pilot Analysis}
\label{chap:a10}

This chapter explains Figure 4.7 (a10\_microscale.png).

% ================================================================
\section{Figure 4.7: Micro-Scale Pilot Analysis}
\label{sec:f4.7}

\subsection{What Type of Figure This Is}
This is a **dual-panel figure**:
\begin{enumerate}
    \item Left: Line graph showing $R_{\mathrm{norm}}$ at very small populations ($N$ = 10 to 1,000)
    \item Right: Bar chart specifically for $N=50$ (village pilot scale)
\end{enumerate}

\subsection{What Is Being Shown}

\textbf{Left Panel (R\_norm vs N):}
\begin{itemize}
    \item X-axis: Population $N$ (10, 25, 50, 75, 100, 250, 500, 1,000)
    \item Y-axis: $R_{\mathrm{norm}}$ (0 to 1.0)
    \item 6 lines, one per scheme
    \item Shows how risk changes at very small populations (village-level pilots)
    \item Vertical dashed line at $N=50$ marks typical VDC pilot size
\end{itemize}

\textbf{Right Panel (N=50):}
\begin{itemize}
    \item X-axis: The six schemes
    \item Y-axis: $R_{\mathrm{norm}}$ at $N=50$
    \item Bars show risk at village pilot scale
    \item Annotations show exact values above each bar
\end{itemize}

Key values at $N=50$:
\begin{itemize}
    \item Plain Hash Token: 0.674
    \item Blind Signature: 0.607
    \item Hybrid Scheme: 0.537
    \item Ring Signature: 0.485
    \item Stealth Address: 0.545
    \item ZKP Token: 0.071
\end{itemize}

\subsection{What This Graph Signifies}
\begin{enumerate}
    \item \textbf{Village pilots (N=50) are the highest-risk scenario} - All schemes except ZKP exceed $R_{\mathrm{norm}} = 0.48$ (nearly half of users identifiable)
    \item \textbf{ZKP Token is the ONLY safe scheme at pilot scale} - At $N=50$, ZKP has $R_{\mathrm{norm}} = 0.071$ while all others are above 0.48
    \item \textbf{The 6.8-fold difference at N=50} is much larger than the 2-fold difference at $N=1,000$ - showing that scheme choice matters most at small scale
    \item \textbf{At N=10, risk approaches maximum} - Plain Hash Token reaches $R_{\mathrm{norm}} = 0.932$, nearly fully identifying all users
    \item \textbf{If NRB plans community pilots, ZKP is essential} - This figure directly supports this policy recommendation
\end{enumerate}

\chapter{Analysis 11: Lambda Sensitivity}
\label{chap:a11}

This chapter explains Figure 4.8 (a11\_lambda\_sensitivity.png).

% ================================================================
\section{Figure 4.8: Lambda Sensitivity Analysis}
\label{sec:f4.8}

\subsection{What Type of Figure This Is}
This is a **dual-panel line graph** showing how $R_{\mathrm{norm}}$ changes as the interaction coefficient $\lambda$ varies from 0 to 100.

\subsection{What Is Being Shown}

\textbf{Left Panel (Plain Hash Token):}
\begin{itemize}
    \item X-axis: $\lambda$ (amount, wallet) from 0 to 100
    \item Y-axis: $R_{\mathrm{norm}}$ (0 to 4.0, capped at 1.0)
    \item Shows how Plain Hash Token's risk grows as interaction strength increases
    \item At very high $\lambda$ (>25), $R_{\mathrm{norm}}$ exceeds 1.0 (theoretical maximum)
\end{itemize}

\textbf{Right Panel (ZKP Token):}
\begin{itemize}
    \item X-axis: $\lambda$ (wallet, timestamp) from 0 to 100
    \item Y-axis: $R_{\mathrm{norm}}$ (flat at 0.035)
    \item Shows a flat horizontal line - ZKP is immune to $\lambda$ changes
\end{itemize}

Key values for Plain Hash Token:
\begin{itemize}
    \item $\lambda=1$: $R_{\mathrm{norm}}=0.202$
    \item $\lambda=5$: $R_{\mathrm{norm}}=0.343$
    \item $\lambda=10$: $R_{\mathrm{norm}}=0.520$
    \item $\lambda=25$: $R_{\mathrm{norm}}=1.049$ (exceeds 1.0)
    \item $\lambda=100$: $R_{\mathrm{norm}}=3.696$
\end{itemize}

\subsection{What This Graph Signifies}
\begin{enumerate}
    \item \textbf{Plain Hash Token is highly sensitive to $\lambda$} - Risk more than doubles when $\lambda$ increases from 1 to 10
    \item \textbf{ZKP Token is completely immune} - The flat line shows $R_{\mathrm{norm}}=0.035$ regardless of $\lambda$
    \item \textbf{At very high interaction ( $\lambda>25$ ), Plain Hash exceeds the maximum} - This is mathematically possible because the theoretical maximum assumes independence
    \item \textbf{Empirical $\lambda$ values for Nepal are near 1} - The model's linear regime applies to real data
    \item \textbf{ZKP's immunity comes from reducing all $s(a_i)$ to near-zero} - With near-zero base values, even high $\lambda$ produces negligible interaction
\end{enumerate}

\chapter{Analysis 12: Worst-Case Scenario}
\label{sec:a12}

This chapter explains Figure 4.9 (a12\_worst\_case.png).

% ================================================================
\section{Figure 4.9: Worst-Case Scenario}
\label{sec:f4.9}

\subsection{What Type of Figure This Is}
This is a **dual-panel figure**:
\begin{enumerate}
    \item Left: Bar chart comparing all schemes to worst-case and theoretical maximum
    \item Right: Grouped bar chart comparing $s(a_i)$ values for Plain Hash Token, ZKP Token, and worst-case
\end{enumerate}

\subsection{What Is Being Shown}

\textbf{Left Panel (R\_norm Comparison):}
\begin{itemize}
    \item X-axis: All 6 schemes + Worst Case + Theoretical Maximum
    \item Y-axis: $R_{\mathrm{norm}}$ (0 to 1.1)
    \item The red dashed line at $R_{\mathrm{norm}}=1.0$ represents the theoretical maximum
    \item Worst Case achieves 0.999 (99.9\% of maximum), validating the normalization
\end{itemize}

\textbf{Right Panel (s(ai) Comparison):}
\begin{itemize}
    \item X-axis: The 7 attributes
    \item Y-axis: $s(a_i)$ values (0 to 1.0)
    \item Three groups of bars: Plain Hash (red), ZKP Token (blue), Worst Case (magenta)
    \item Shows how each attribute's uniqueness differs across configurations
\end{itemize}

Worst-case configuration:
\begin{itemize}
    \item $s(\text{amount}) = 0.984$
    \item $s(\text{timestamp}) = 1.000$
    \item $s(\text{merchant}) = 1.000$
    \item $s(\text{wallet}) = 1.000$
    \item $s(\text{frequency}) = 1.000$
    \item $s(\text{payment\_channel}) = 1.000$
    \item $s(\text{transaction\_zone}) = 1.000$
    \item $R_{\mathrm{raw}} = 27.984$
    \item $R_{\mathrm{norm}} = 0.999$ (99.9\% of maximum)
\end{itemize}

\subsection{What This Graph Signifies}
\begin{enumerate}
    \item \textbf{The normalization is correctly calibrated} - Worst case achieves 99.9\% of the theoretical maximum, confirming $R_{\mathrm{max}}=28$ is not an overestimate
    \item \textbf{All deployed schemes are far below the ceiling} - Even Plain Hash at 0.317 is only 31.7\% of maximum
    \item \textbf{ZKP achieves 3.5\% of worst case} - This represents the state-of-the-art in CBDC privacy
    \item \textbf{Plain Hash is not acceptable despite being below ceiling} - 31.7\% of maximum still means 31.7\% of users are structurally vulnerable
    \item \textbf{The gap of 0.001 from maximum} exists because amount follows an exponential distribution (not perfectly distinct), leaving a tiny fraction sharing values
\end{enumerate}

\chapter{Summary}
\label{chap:summary}

This tutorial has covered all figures in the thesis:

\begin{enumerate}
    \item \textbf{Figure 4.1a}: Bar chart showing $R_{\mathrm{norm}}$ ranking of six cryptographic schemes
    \item \textbf{Figure 4.1b}: Stacked bar showing attribute contributions to risk
    \item \textbf{Figure 4.1c}: Nepal-specific attributes showing zero contribution (due to demographic clustering)
    \item \textbf{Figure 4.2}: Heatmaps showing $\lambda_{ij}$ interaction coefficients
    \item \textbf{Figure 4.3}: Population scalability (N = 100 to 100,000)
    \item \textbf{Figure 4.4}: Mitigation effectiveness (91-97\% risk reduction)
    \item \textbf{Figure 4.5}: Attribute sensitivity (m = 1 to 7)
    \item \textbf{Figure 4.6}: Seed sensitivity validation (30 runs)
    \item \textbf{Figure 4.7}: Micro-scale pilot analysis (N = 10 to 1,000)
    \item \textbf{Figure 4.8}: Lambda sensitivity ( $\lambda$ = 0 to 100)
    \item \textbf{Figure 4.9}: Worst-case calibration (99.9\% of maximum)
\end{enumerate}

Key takeaways:
\begin{itemize}
    \item \textbf{Wallet is the dominant risk driver} across all schemes
    \item \textbf{ZKP Token provides the strongest privacy} (lowest $R_{\mathrm{norm}}$) but has highest computational cost
    \item \textbf{Nepal's demographic structure provides passive protection} on payment\_channel and transaction\_zone
    \item \textbf{Scale cannot substitute for cryptographic design} - Plain Hash remains risky even at N=100,000
    \item \textbf{Mitigation strategies are effective} - wallet unlinkability is the single most impactful
    \item \textbf{Village pilots (N=50) require ZKP} - all other schemes exceed 48\% risk at this scale
\end{itemize}

\end{document}